% =============================================================================
\chapter{Gradients de Politique --- REINFORCE}
\label{ch:gradient_politique}
% =============================================================================

Jusqu'ici, nous avons appris des \emph{fonctions de valeur} et d\'eduit la politique indirectement. Les m\'ethodes de gradient de politique inversent cette logique~: elles param\'etrisent directement la politique $\pi_\theta$ et optimisent ses param\`etres par mont\'ee de gradient. L'algorithme REINFORCE, propos\'e par Ronald Williams en 1992, utilise le \emph{th\'eor\`eme du gradient de politique} --- un r\'esultat \'el\'egant qui exprime le gradient de la performance esp\'er\'ee en termes d'esp\'erance sous la politique, sans n\'ecessiter de diff\'erencier \`a travers la dynamique de l'environnement. Cette approche ouvre la voie aux espaces d'actions continus et aux politiques stochastiques expressives.

\begin{intuition}
Les m\'ethodes de gradient de politique optimisent directement les param\`etres
$\theta$ d'une politique $\pi_\theta$ en montant le gradient de la performance
$J(\theta) = \E_{\pi_\theta}[G_0]$. Contrairement aux m\'ethodes fond\'ees
sur la valeur, elles peuvent naturellement g\'erer les espaces d'action continus
et les politiques stochastiques.
\end{intuition}

% -----------------------------------------------------------------------------
\section{Motivation}
% -----------------------------------------------------------------------------

Les m\'ethodes de type DQN pr\'esentent des limitations~:
\begin{itemize}
    \item Elles ne s'appliquent qu'aux espaces d'action \textbf{discrets} et de
    petite taille (le $\max$ sur les actions est explicite).
    \item La politique est implicite (gloutonne par rapport \`a $Q$) et ne peut
    repr\'esenter que des politiques d\'eterministes.
    \item De petits changements dans $Q$ peuvent causer de grands changements
    dans la politique, rendant l'apprentissage instable.
\end{itemize}

L'id\'ee du gradient de politique est de param\'etrer directement $\pi_\theta$
et d'optimiser $\theta$ par mont\'ee de gradient.

% -----------------------------------------------------------------------------
\section{Objectif et th\'eor\`eme du gradient de politique}
% -----------------------------------------------------------------------------

\begin{definition}[Objectif de performance]
L'objectif de performance est le retour esp\'er\'e sous la politique $\pi_\theta$~:
\[
    J(\theta) = \E_{\tau \sim \pi_\theta}\left[\sum_{t=0}^{T-1} \gamma^t R_{t+1}\right]
    = \E_{s_0 \sim \rho_0}\left[V^{\pi_\theta}(s_0)\right]
\]
o\`u $\rho_0$ est la distribution initiale des \'etats et $\tau = (s_0, a_0, r_1, s_1, \ldots)$
est une trajectoire.
\end{definition}

\begin{theoreme}[Th\'eor\`eme du gradient de politique]
\label{thm:policy_gradient}
Le gradient de $J(\theta)$ est~:
\[
    \nabla_\theta J(\theta) = \E_{\pi_\theta}\left[
    \sum_{t=0}^{T-1} \nabla_\theta \log \pi_\theta(A_t \mid S_t) \, Q^{\pi_\theta}(S_t, A_t)
    \right].
\]
\end{theoreme}

\begin{proof}
On commence par le cas \`a un pas. Soit $J(\theta) = \sum_s d^{\pi_\theta}(s)
\sum_a \pi_\theta(a|s) Q^{\pi_\theta}(s,a)$ o\`u $d^{\pi_\theta}(s)$ est la
distribution stationnaire. Alors~:
\begin{align*}
    \nabla_\theta J(\theta)
    &= \sum_s d^{\pi_\theta}(s) \sum_a \nabla_\theta \pi_\theta(a|s) Q^{\pi_\theta}(s,a) \\
    &\quad + \sum_s \nabla_\theta d^{\pi_\theta}(s) \sum_a \pi_\theta(a|s) Q^{\pi_\theta}(s,a).
\end{align*}
Le r\'esultat remarquable (Sutton et al., 2000) est que l'on peut ignorer le
second terme (d\'ependance de $d^{\pi_\theta}$ en $\theta$). En utilisant
l'identit\'e du log-gradient $\nabla_\theta \pi_\theta = \pi_\theta \nabla_\theta \log \pi_\theta$~:
\[
    \nabla_\theta J(\theta) = \E_{\pi_\theta}\left[
    \nabla_\theta \log \pi_\theta(A_t | S_t) \, Q^{\pi_\theta}(S_t, A_t)\right].
\]
Pour l'horizon fini, on somme sur $t = 0, \ldots, T-1$.
\end{proof}

\begin{formules}[title=\textbf{Th\'eor\`eme du gradient de politique}]
\[
    \nabla_\theta J(\theta) = \E_{\pi_\theta}\left[
    \sum_{t=0}^{T-1} \nabla_\theta \log \pi_\theta(A_t | S_t) \, Q^{\pi_\theta}(S_t, A_t)
    \right]
\]
\end{formules}

% -----------------------------------------------------------------------------
\section{Algorithme REINFORCE}
% -----------------------------------------------------------------------------

\begin{definition}[REINFORCE]
L'algorithme \textbf{REINFORCE} (Williams, 1992) estime le gradient de politique
en rempla\c{c}ant $Q^{\pi_\theta}(S_t, A_t)$ par le retour Monte Carlo $G_t$~:
\[
    \nabla_\theta J(\theta) \approx \frac{1}{N} \sum_{i=1}^{N} \sum_{t=0}^{T_i-1}
    \nabla_\theta \log \pi_\theta(a_t^{(i)} | s_t^{(i)}) \, G_t^{(i)}.
\]
\end{definition}

\begin{algorithme}[title=\textbf{REINFORCE}]
\begin{enumerate}
    \item Initialiser les param\`etres $\theta$ de $\pi_\theta$
    \item Pour chaque \'episode~:
    \begin{enumerate}
        \item G\'en\'erer une trajectoire $\tau = (s_0, a_0, r_1, \ldots, s_T)$
        en suivant $\pi_\theta$
        \item Pour $t = T-1, T-2, \ldots, 0$~: calculer $G_t = \sum_{k=0}^{T-t-1} \gamma^k r_{t+k+1}$
        \item $\theta \leftarrow \theta + \alpha \sum_{t=0}^{T-1} \gamma^t \nabla_\theta \log \pi_\theta(a_t | s_t) \, G_t$
    \end{enumerate}
\end{enumerate}
\end{algorithme}

% -----------------------------------------------------------------------------
\section{R\'eduction de variance~: la ligne de base}
% -----------------------------------------------------------------------------

\begin{theoreme}[Ligne de base]
\label{thm:baseline}
Pour toute fonction $b(s)$ ne d\'ependant pas de $a$~:
\[
    \E_{\pi_\theta}\left[\nabla_\theta \log \pi_\theta(a|s) \, b(s)\right] = 0.
\]
On peut donc soustraire $b(s)$ du retour sans introduire de biais~:
\[
    \nabla_\theta J(\theta) = \E_{\pi_\theta}\left[
    \nabla_\theta \log \pi_\theta(a_t|s_t) \, (G_t - b(s_t))
    \right].
\]
\end{theoreme}

\begin{proof}
\begin{align*}
    \E_\pi[\nabla_\theta \log \pi_\theta(a|s) \, b(s)]
    &= \sum_a \pi_\theta(a|s) \frac{\nabla_\theta \pi_\theta(a|s)}{\pi_\theta(a|s)} b(s) \\
    &= b(s) \sum_a \nabla_\theta \pi_\theta(a|s) \\
    &= b(s) \nabla_\theta \underbrace{\sum_a \pi_\theta(a|s)}_{=1} = 0. \qedhere
\end{align*}
\end{proof}

\begin{remarque}
Le choix optimal de la ligne de base pour minimiser la variance est~:
\[
    b^*(s) = \frac{\E[\norm{\nabla_\theta \log \pi_\theta}^2 G_t \mid S_t = s]}
    {\E[\norm{\nabla_\theta \log \pi_\theta}^2 \mid S_t = s]}.
\]
En pratique, on utilise $b(s) \approx V^{\pi_\theta}(s)$, ce qui donne le
gradient \`a base d'avantage~:
\[
    \nabla_\theta J(\theta) \approx \E_{\pi_\theta}\left[
    \nabla_\theta \log \pi_\theta(a_t|s_t) \, \hat{A}_t
    \right]
\]
o\`u $\hat{A}_t = G_t - V(s_t; \phi)$ est l'estimateur de l'avantage.
C'est la base des m\'ethodes acteur-critique (Chapitre~\ref{ch:acteur_critique}).
\end{remarque}

% -----------------------------------------------------------------------------
\section{Politique gaussienne pour actions continues}
% -----------------------------------------------------------------------------

\begin{definition}[Politique gaussienne]
Pour un espace d'action continu $\mathcal{A} \subseteq \R^d$, on param\`etre~:
\[
    \pi_\theta(a | s) = \mathcal{N}\!\left(a; \mu_\theta(s), \sigma_\theta^2(s) I\right)
\]
o\`u $\mu_\theta(s)$ et $\sigma_\theta(s)$ sont des sorties du r\'eseau de neurones.
Le log-gradient est~:
\[
    \nabla_\theta \log \pi_\theta(a|s) =
    \frac{(a - \mu_\theta(s))}{\sigma_\theta^2(s)} \nabla_\theta \mu_\theta(s)
    + \left(\frac{(a - \mu_\theta(s))^2}{\sigma_\theta^3(s)} - \frac{1}{\sigma_\theta(s)}\right) \nabla_\theta \sigma_\theta(s).
\]
\end{definition}

% -----------------------------------------------------------------------------
\section{Politique softmax pour actions discr\`etes}
% -----------------------------------------------------------------------------

\begin{definition}[Politique softmax]
Pour un espace d'action discret $\mathcal{A} = \{1, \ldots, K\}$~:
\[
    \pi_\theta(a | s) = \frac{\exp(h_\theta(s, a))}{\sum_{a'} \exp(h_\theta(s, a'))}
\]
o\`u $h_\theta(s, a)$ sont les logits produits par le r\'eseau de neurones.
Le log-gradient est~:
\[
    \nabla_\theta \log \pi_\theta(a|s) = \nabla_\theta h_\theta(s,a)
    - \sum_{a'} \pi_\theta(a'|s) \nabla_\theta h_\theta(s,a').
\]
\end{definition}

% -----------------------------------------------------------------------------
\section{Causalit\'e et r\'eduction de variance suppl\'ementaire}
% -----------------------------------------------------------------------------

\begin{proposition}[Causalit\'e dans le gradient de politique]
Les r\'ecompenses pass\'ees ne d\'ependent pas des actions futures. On peut
donc remplacer $G_t$ par le \textbf{retour futur}~:
\[
    \nabla_\theta J(\theta) = \E_{\pi_\theta}\left[
    \sum_{t=0}^{T-1} \nabla_\theta \log \pi_\theta(a_t|s_t)
    \left(\sum_{k=t}^{T-1} \gamma^{k-t} r_{k+1}\right)
    \right].
\]
Cela r\'eduit la variance sans introduire de biais.
\end{proposition}

% -----------------------------------------------------------------------------
\section{Impl\'ementation PyTorch}
% -----------------------------------------------------------------------------

\begin{codeblock}[title=\textbf{REINFORCE avec ligne de base}]
\begin{minted}{python}
import torch
import torch.nn as nn
import torch.optim as optim
from torch.distributions import Categorical
import gymnasium as gym
import numpy as np

class PolicyNetwork(nn.Module):
    def __init__(self, obs_dim, n_actions, hidden=128):
        super().__init__()
        self.net = nn.Sequential(
            nn.Linear(obs_dim, hidden), nn.ReLU(),
            nn.Linear(hidden, hidden), nn.ReLU(),
            nn.Linear(hidden, n_actions)
        )

    def forward(self, x):
        logits = self.net(x)
        return Categorical(logits=logits)

class ValueNetwork(nn.Module):
    def __init__(self, obs_dim, hidden=128):
        super().__init__()
        self.net = nn.Sequential(
            nn.Linear(obs_dim, hidden), nn.ReLU(),
            nn.Linear(hidden, hidden), nn.ReLU(),
            nn.Linear(hidden, 1)
        )

    def forward(self, x):
        return self.net(x).squeeze(-1)

def reinforce_baseline(env_name="CartPole-v1", n_episodes=1000,
                       gamma=0.99, lr_pi=1e-3, lr_v=1e-3):
    env = gym.make(env_name)
    obs_dim = env.observation_space.shape[0]
    n_actions = env.action_space.n

    policy = PolicyNetwork(obs_dim, n_actions)
    value_fn = ValueNetwork(obs_dim)
    opt_pi = optim.Adam(policy.parameters(), lr=lr_pi)
    opt_v = optim.Adam(value_fn.parameters(), lr=lr_v)

    for ep in range(n_episodes):
        states, actions, rewards = [], [], []
        s, _ = env.reset()
        done = False

        while not done:
            s_t = torch.FloatTensor(s)
            dist = policy(s_t)
            a = dist.sample()
            s_next, r, terminated, truncated, _ = env.step(a.item())
            states.append(s)
            actions.append(a.item())
            rewards.append(r)
            s = s_next
            done = terminated or truncated

        # Calcul des retours
        T = len(rewards)
        returns = np.zeros(T)
        G = 0
        for t in reversed(range(T)):
            G = rewards[t] + gamma * G
            returns[t] = G

        states_t = torch.FloatTensor(np.array(states))
        actions_t = torch.LongTensor(actions)
        returns_t = torch.FloatTensor(returns)

        # Mise a jour de la ligne de base (value function)
        values = value_fn(states_t)
        value_loss = nn.MSELoss()(values, returns_t)
        opt_v.zero_grad()
        value_loss.backward()
        opt_v.step()

        # Mise a jour de la politique
        with torch.no_grad():
            advantages = returns_t - value_fn(states_t)

        dist = policy(states_t)
        log_probs = dist.log_prob(actions_t)
        policy_loss = -(log_probs * advantages).mean()
        opt_pi.zero_grad()
        policy_loss.backward()
        opt_pi.step()

        if ep % 100 == 0:
            print(f"Episode {ep}, Return: {returns[0]:.1f}")

    return policy

policy = reinforce_baseline()
\end{minted}
\end{codeblock}

\begin{codeblock}[title=\textbf{REINFORCE pour actions continues}]
\begin{minted}{python}
from torch.distributions import Normal

class GaussianPolicy(nn.Module):
    def __init__(self, obs_dim, act_dim, hidden=64):
        super().__init__()
        self.shared = nn.Sequential(
            nn.Linear(obs_dim, hidden), nn.Tanh(),
            nn.Linear(hidden, hidden), nn.Tanh()
        )
        self.mu_head = nn.Linear(hidden, act_dim)
        self.log_std = nn.Parameter(torch.zeros(act_dim))

    def forward(self, x):
        h = self.shared(x)
        mu = self.mu_head(h)
        std = self.log_std.exp()
        return Normal(mu, std)
\end{minted}
\end{codeblock}

% -----------------------------------------------------------------------------
\section{Analyse de la variance}
% -----------------------------------------------------------------------------

\begin{theoreme}[Variance de REINFORCE]
La variance de l'estimateur REINFORCE cro\^it avec~:
\begin{enumerate}
    \item La longueur de l'horizon $T$~: $\Var[\hat{g}] = O(T)$.
    \item La dimensionnalit\'e de $\theta$~: chaque composante a sa propre variance.
    \item La stochasticit\'e de la politique et de l'environnement.
\end{enumerate}
La ligne de base r\'eduit la variance d'un facteur qui d\'epend de la corr\'elation
entre $\nabla \log \pi$ et $G_t - b(s)$.
\end{theoreme}

\begin{attention}[title=\textbf{Limitations de REINFORCE}]
REINFORCE souffre de~:
\begin{itemize}
    \item Haute variance $\Rightarrow$ convergence lente.
    \item N\'ecessite des \'episodes complets (pas de bootstrap).
    \item Utilisation inefficace des donn\'ees (on-policy, pas de rejeu).
\end{itemize}
Ces limitations motivent les m\'ethodes acteur-critique (Chapitre~\ref{ch:acteur_critique}).
\end{attention}

% -----------------------------------------------------------------------------
\section{Exercices}
% -----------------------------------------------------------------------------

\begin{exercice}[Effet de la ligne de base]
Impl\'ementez REINFORCE avec et sans ligne de base sur \texttt{CartPole-v1}.
Tracez la variance du gradient estim\'e au cours de l'entra\^inement pour les
deux versions. Comparez la vitesse de convergence.
\end{exercice}

\begin{exercice}[Actions continues]
Appliquez REINFORCE avec politique gaussienne \`a l'environnement
\texttt{Pendulum-v1} (action continue). Utilisez un r\'eseau pour $\mu$ et
un param\`etre apprenable pour $\log \sigma$.
\end{exercice}

\begin{exercice}[Preuve du th\'eor\`eme]
D\'emontrez le th\'eor\`eme du gradient de politique pour le cas \`a horizon
infini en utilisant la distribution stationnaire $d^{\pi_\theta}$.
\end{exercice}

\begin{exercice}[Ligne de base optimale]
D\'erivez analytiquement la ligne de base $b^*(s)$ qui minimise la variance
de l'estimateur du gradient et montrez que $V^{\pi_\theta}(s)$ en est une
bonne approximation.
\end{exercice}
